\TieudeT{PHONG TỎA VẬT LÍ}
\subsubsection*{PHẦN I. Thí sinh trả lời câu hỏi từ câu 1 đến câu 18. Mỗi câu hỏi thí sinh chỉ chọn một phương án}
\setcounter{ex}{0}
%%mac
%%\loai{Nhận biết phương trình}
%\begin{ex}%book
%	Tọa độ một chất điểm chuyển động thẳng đều có dạng nào sau đây?
%	\choice
%	{\True $x-x_0=v t$}
%	{$x=v+x_0 t$}
%	{$x=\left(x_0+v\right) t$}
%	{$x+x_0=v t$}
%	\loigiai{
%	}
%	\haidong{7} \end{ex}
%\begin{ex}%book
%	Phương trình nào sau đây mô tả chuyển động thẳng đều?
%	\choice
%	{$x=2\sin t$}
%	{$x=2+3 t^2$}
%	{$x=-2 t^2+6$}
%	{\True $x=-4 t+3$}
%	\loigiai{
%	}
%	\haidong{7} \end{ex}
%%\loai{Xác định các đại lượng từ phương trình}
%\begin{ex}%book
%	Một vật chuyển động thẳng đều có phương trình $x=10+2t$ ($x$ đo bằng $m$, $t$ đo bằng $s$). Vận tốc chuyển động của vật là
%	\choice
%	{$10\ m/s$}
%	{\True $2\ m/s$}
%	{$12\ m/s$}
%	{$5\ m/s$}
%	\loigiai{
%	}
%	\haidong{7} \end{ex}
%\begin{ex}%book
%	Phương trình chuyển động của một chất điểm dọc theo theo trục $Ox$ có dạng $x=5+12 t$ ($x$ đo bằng $km$, $t$ đo bằng $h$). Chất điểm đó xuất phát
%	\choice
%	{từ điểm cách $O$ là $5\ km$, với vận tốc $60\ km/h$}
%	{từ điểm $O$, với vận tốc $12\ km/h$}
%	{từ điểm $O$, với vận tốc $60\ km/h$}
%	{\True từ điểm cách $O$ là $5\ km$, với vận tốc $12\ km / h$}
%	\haidong{3}
%	\loigiai{
%	}
%	\haidong{7} \end{ex}	
%
%%\loai{Viết phương trình chuyển động của một xe}
%\begin{ex}%book %[1P1B3-1][Loai1] 
%	Một ô tô chuyển động trên một đoạn đường thẳng và có tốc độ luôn luôn bằng $80\  km/h$. Bến xe nằm ở đầu đoạn đường và xe ô tô xuất phát từ một địa điểm cách bến xe $3\ km$. Chọn bến xe làm vật mốc, chọn thời điểm ô tô xuất phát làm mốc thời gian, chiều dương là chiều chuyển động của xe. Phương trình chuyển động của xe ô tô trên đoạn đường thẳng này là
%	\choice
%	{\True $x = 3 + 80t\ (km)$}
%	{$x = (80 - 3)t\ (km)$}
%	{$x = 3 - 80t\ (km)$}
%	{$x = 80t\ (km)$}
%	\haidong{6}
%	\loigiai{
%		Chọn gốc tọa độ tại bến xe, ta có:
%		$x_0=3\ km;\ v=80\ km/h\Rightarrow x = x_0 + vt = 3 + 80t$ (km).
%	}
%	\haidong{7} \end{ex}
%\begin{ex}%book%[1P1K3-1][Loai1]
%	Một xe taxi chuyển động trên một đoạn đường thẳng và có tốc độ không đổi là $45\ km/h$. Bến xe nằm ở đầu đường thẳng và taxi xuất phát từ một điểm cách bến xe $5\ km$. Chọn bến xe làm vật mốc, chọn thời điểm taxi xuất phát làm mốc thời gian và chọn chiều chuyển động của taxi làm chiều dương. Phương trình chuyển động của taxi trên đoạn đường thẳng này là
%	\choice
%	{\True $x = 5 + 45t$}
%	{$x = 45 - 5t$}
%	{$x = 5 - 45t$}
%	{$x = 45t$}
%	\haidong{4}
%	\loigiai{
%		\begin{itemize}
%			\item Mốc là bến xe  $\Rightarrow x_0=5\ km$.
%			\item Chiều dương là chiều chuyển động của xe $\Rightarrow v=45\ km/h$. 
%			\item Phương trình chuyển động của xe là: $x=5+45t$. 
%		\end{itemize}
%	}
%	%<MyLT>
%	\haidong{7} \end{ex}
%%\loai{Viết phương trình chuyển động của hai xe}
%\begin{ex}%book%[1P1B3-1][Loai1]
%	Cùng một lúc tại hai điểm $A$ và $B$ cách nhau $15\ km$ có hai ô tô chạy theo chiều $AB$. Tốc độ của ô tô chạy từ $A$ là $40\ km/h$ và của ô tô chạy từ $B$ là $35\ km/h$. Chọn $A$ làm mốc, thời điểm xuất phát của hai xe ô tô làm mốc thời gian và chiều chuyển động của hai xe làm chiều dương. Chọn gốc thời gian lúc vật bắt đầu chuyển động. Phương trình chuyển động của các ô tô chạy từ $A$ và từ $B$ lần lượt là
%	\choice
%	{\True $x_A = 40t;\ x_B = 35t + 15$}
%	{$x_A = 40t + 15;\ x_B = 35t$}
%	{$x_A = 40t;\ x_B = 35t - 15$}
%	{$x_A = -40t;\ x_B = 35t$}
%	\haidong{4}
%	\loigiai{
%		\begin{itemize}
%			\item A là mốc $\Rightarrow$ tọa độ ban đầu của 2 xe là: $x_{0A}=0\ km;\ x_{0B}=15\ km$. 
%			\item Hai xe chuyển động theo chiều dương: $\Rightarrow v_A=40\ km;\ v_B=35\ km$.
%			\item Phương trình chuyển động của vật là: $x=x_0+vt$.
%			\item Phương trình chuyển động của ô tô chạy từ A và B lần lượt là: $x_A=40t;\ x_B=15+35t$.
%		\end{itemize}
%	}
%	\haidong{7} \end{ex}

\begin{ex}%book
Nếu đồ thị độ dịch chuyển $(d)$ - thời gian $(t)$ của một vật là một đoạn thẳng dốc xuống thì ta có thể biết được
\choice
{độ dịch chuyển âm}
{vật chuyển động theo chiều dương của trục tọa độ}
{độ dịch chuyển dương}
{\True vật chuyển động ngược chiều dương của trục tọa độ}

\loigiai{
}
\haidong{7} \end{ex}
\begin{ex}%book
Đồ thị độ dịch chuyển - thời gian khi chọn gốc tọa độ $O$ là điểm xuất phát, chiều dương trùng chiều chuyển động, mốc thời gian tính từ lúc xuất phát thì đồ thị sẽ có dạng
\choice
{là đường thẳng không đi qua $O$ và xiên lên}
{là đường thẳng đi qua $O$ và xiên xuống}
{\True là đường thẳng đi qua $O$ và xiên lên}
{là đường thẳng không đi qua $O$ và xiên lên}
\loigiai{
}
\haidong{7} %<MyLT1>
\end{ex}
\begin{ex}%book%7.2
\immini[thm]{Theo đồ thị ở hình bên, vật chuyển động thẳng đều trong khoảng thời gian
\choice
{từ 0 đến $t_2$}
{từ $t_1$ đến $t_2$}
{\True từ 0 đến $t_1$ và từ $t_2$ đến $t_3$}
{từ 0 đến $t_3$}}		{\begin{tikzpicture}[scale=.7,thick,>=stealth',draw=Mapcolor]
\tkzDefPoints{0/0/o, 0/0/i, 7.5/0/t, 0/4/d}
\tkzDrawSegments[thick,Mapcolor,->](o,t i,d)
\tkzLabelPoint[left](o){$O$}
\node[right] at (d){d};
\node[above] at (t){t};
\draw[blue,line width=1.2pt](0,0)--(2,3)--(5,3)--(7,0)node[below,black]{$t_3$};
\draw[dashed](2,3)--(2,0)node[below]{$t_1$};
\draw[dashed](5,3)--(5,0)node[below]{$t_2$};
\end{tikzpicture}}
\haidong{7} \end{ex}
\begin{ex}%book
Trong các đồ thị tọa độ - thời gian $(x-t)$ dưới đây, đồ thị nào \textbf{không} biểu diễn chuyển động thẳng đều.
\begin{center}
\begin{tikzpicture}[scale=1,thick,>=stealth',x=1.5cm,y=0.2cm,draw=Mapcolor]
\tkzDefPoints{0/0/o, 0/12/x, 2/0/t, 0/4/m, 1.2/10/n, 2/0/n1, 0/12/n2}
\tkzDrawLines[very thick,color=Mapcolor,add = 0 and 0](m,n)
\tkzDrawSegments[->](o,x o,t)
\tkzLabelPoint[below](o){$O$}
\tkzLabelPoint[right](x){$x$}
\tkzLabelPoint[above](t){$t\ (s)$}
\tkzLabelPoint[left](m){$x_0$}
\node at (1,-3.5) {Hình A}; % Thêm nhãn cho đồ thị 1
\begin{scope}[shift={(3,0)}]
\tkzDefPoints{0/5/o', 0/0/o, 0/12/x, 2/0/t, 0/5/m, 1.9/5/n, 2/0/n1, 0/12/n2}
\tkzDrawLines[very thick,color=Mapcolor,add = 0 and 0](m,n)
\tkzDrawSegments[->](o,x o,t)
\tkzLabelPoint[right](x){$x$}
\tkzLabelPoint[above](t){$t\ (s)$}
\tkzLabelPoint[below left](o){$O$}
\tkzLabelPoint[left](m){$x_0$}
\node at (1,-3.5) {Hình B}; % Thêm nhãn cho đồ thị 2
\end{scope}
\begin{scope}[shift={(6,0)}]
\tkzDefPoints{0/0/o, 0/12/x, 2/0/t, 0/5/m, 1/0/n, 2/0/n1, 0/12/n2}
\tkzDrawLines[very thick,color=Mapcolor,add = 0 and 0](m,n)
\tkzDrawSegments[->](o,x o,t)
\tkzLabelPoint[below](o){$O$}
\tkzLabelPoint[right](x){$x$}
\tkzLabelPoint[above](t){$t\ (s)$}
\tkzLabelPoint[left](m){$x_0$}
\node at (1,-3.5) {Hình C}; % Thêm nhãn cho đồ thị 2
\end{scope}
\begin{scope}[shift={(9,0)}]
\tkzDefPoints{0/5/o', 0/0/o, 0/12/x, 2/5/t, 0/2/m, 1.4/10/n, 2/0/n1, 0/12/n2}
\tkzDrawLines[very thick,color=Mapcolor,add = 0 and 0](m,n)
\tkzDrawSegments[->](o,x o',t)
\tkzLabelPoint[left](o'){$O$}
\tkzLabelPoint[right](x){$x$}
\tkzLabelPoint[above](t){$t\ (s)$}
\tkzLabelPoint[left](m){$x_0$}
\node at (1,-3.5) {Hình D}; % Thêm nhãn cho đồ thị 2
\end{scope}

\end{tikzpicture}
\end{center}

\choice
{Đồ thị hình $A$}
{\True Đồ thị hình $B$}
{Đồ thị hình $C$}
{Đồ thị hình $D$}
\loigiai{
}
\haidong{7} \end{ex}
\begin{ex}%book
\immini[thm]{Cho đồ thị độ dịch chuyển - thời gian của một vật như hình. Trong $3$ giây đầu vật dịch chuyển một đoạn bao nhiêu?
\choice[2]
{\True $15\ m$}
{$20\ m$}
{$5\ m$}
{$35\ m$}}{\begin{hinhanh}{6}\begin{tikzpicture}[very thick,draw=Mapcolor]
\tkzDefPoints{0/0/o, 0/3/d, 4/0/t}
\tkzDrawSegments[->, thick](o,d o,t)
\draw [dashed, thin](0,1.8) -- (1.8,1.8) -- (1.8,0);
\draw[dashed, thin] (0,2.5) -- (2.5,2.5) -- (2.5,0);
\draw[Mapcolor, very thick] (0,0) -- (2.8,2.8);
\tkzLabelPoint[below](o){O}
\tkzLabelPoint[above right](d){$d ~(m)$}
\tkzLabelPoint[below](t){$t~(s)$}

\node[left] at (0,1.8) {15};
\node[left] at (0,2.5) {20};
\node [below]at (1.8,0) {3};
\node[below] at (2.5,0) {4};
\end{tikzpicture}\end{hinhanh}
}
\loigiai{
}
\haidong{7} \end{ex}
\begin{ex}%book
\immini[thm]{Cho một vật chuyển động tuân theo đồ thị độ dịch chuyển - thời gian như hình. Tính độ dịch chuyển từ lúc $t_1=4$ giờ đến lúc $t_2=7$ giờ.
\choice[2]
{$4\ km$}
{$-6\ km$}
{\True $-4\ km$}
{$6\ km$}}{\begin{hinhanh}{6}\begin{tikzpicture}[very thick,draw=Mapcolor]
\tkzDefPoints{0/0/o, 0/3/d, 4.5/0/t}
\tkzDrawSegments[->, thick](o,d o,t)
\draw [Mapcolor, very thick](0,0) -- (1,2.5)  -- (3,2.5) -- (4,0);
\draw[dashed, thin] (0,2.5) -- (1,2.5) -- (1,0);
\draw[dashed, thin] (3,2.5) -- (3,0);
\node[left] at (0,2.5) {4};
\node [below]at (1,0) {2};
\node [below]at (3,0) {5};
\node[below] at (4,0) {7};
\tkzLabelPoint[below left](o){O}
\tkzLabelPoint[above](d){$d~(km)$}
\tkzLabelPoint[above](t){$t~(h)$}
\end{tikzpicture}\end{hinhanh}}
\haidong{3}
\loigiai{
}
\haidong{7} \end{ex}
\begin{ex}%book
\immini[thm]{Hình vẽ dưới đây mô tả độ dịch chuyển - thời gian của một chiếc xe ô tô chạy trên một đường thẳng. Vận tốc trung bình của xe có giá trị
\choice[2]
{$-90\ km/h$}
{$90\ km/h$}
{$-45\ km/h$}
{\True $45\ km/h$}}{\begin{hinhanh}{6}\begin{tikzpicture}[very thick,draw=Mapcolor]
\tkzDefPoints{0/0/o, 0/3/d, 4/0/t}
\tkzDrawSegments[->, thick](o,d o,t)
\draw  [help  lines, black, dashed, xstep=0.7cm,ystep=0.5cm]  (0,0) node (v1) {}  grid  (3.5,2.5);
\tkzLabelPoint[below left](o){O}
\tkzLabelPoint[above](d){$d~(m)$}
\tkzLabelPoint[above right](t){$t~(s)$}
\draw[Mapcolor] (0,0) -- (3.5,2.5);
\node [left]at (0,1) {45};
\node [left]at (0,2) {90};
\node[below] at (1.4,0) {1};
\node [below]at (2.8,0) {2};
\end{tikzpicture}\end{hinhanh}}
\haidong{3}
\loigiai{
}
\haidong{7} \end{ex}
\begin{ex}%Tailieuchuan
\immini[thm]{Một chất điểm chuyển động có quỹ đạo là đường nét đứt như hình vẽ. Hình vẽ vectơ vận tốc của chất điểm tại điểm nào là đúng?
\choice
{\True Điểm A}
{Điểm C}
{Điểm E}
{Điểm C và E}}{
\begin{tikzpicture}[scale=1,thick,>=stealth',draw=Mapcolor]
\path
(0,0)coordinate(O)
(1.3,1.5)coordinate(A) 
(2.9,1)coordinate(C) 
(6,1)coordinate(E)
;
\draw[dashed,Mapcolor] (0,0) .. controls (0.7,0.5) and (1.5,2) .. (2,2) .. controls (2.6,1.95) and (3.45,-0.5) .. (4,-0.5) .. controls (4.5,-0.5) and (5.4,1) .. (6,1) .. controls (6.65,1) and (7,0.5) .. (7.5,0);

\draw[->,Mapcolor] (1.3,1.5) -- (2.1,2.3);
\draw[->,Mapcolor] (2.9,1) -- (3.7,0.6);
\draw[->,Mapcolor] (6,1) -- (5.6,0.4);
\foreach \x/\g in {A/150,C/60,E/90}
\draw[fill=white](\x) circle(0.05)+(\g:0.3)node{$\x$};
\end{tikzpicture}}
\loigiai{

}
\end{ex}
\begin{ex}%Tailieuchuan
\impicinpar{Một ô tô chạy thử nghiệm trên một đoạn đường thẳng. Cứ $5\ s$ thì có một giọt dầu từ động cơ của ô tô rơi thẳng xuống mặt đường. Hình bên cho thấy mô hình các giọt dầu để lại trên mặt đường. Ô tô chuyển động trên đường này với tốc độ trung bình là
\choice
{$12,5 \ m/s$}
{\True $15 \ m/s$}
{$30 \ m/s$}
{$25 \ m/s$}}{\includegraphics[width=4cm]{images/2025_04_24_b6f4a66b1170539a4b52g-16(2)}}
\loigiai{

}
\end{ex}
\begin{ex}%Tailieuchuan
\immini[thm]{Phát biểu nào đúng về vận tốc của chất điểm trên quỹ đạo (đường nét đứt như hình) tại hai điểm $A$ và $C$?
\choice
{Tại $A$ vật chuyển động nhanh hơn tại $C$}
{Tại $A$ và tại $C$ vật có cùng tốc độ}
{Tại $A$ và tại $C$ vật đang chuyển động cùng hướng}
{\True Tại $A$ vật chuyển động chậm hơn tại $C$}}{\begin{hinhanh}{6}\begin{tikzpicture}[scale=1,thick,>=stealth',draw=Mapcolor]
\path
(0,0)coordinate(O)
(0.2,0.1)coordinate(A) 
(2,2)coordinate(C) 
;
\draw[dashed,Mapcolor] (0,0) .. controls (0.7,0.2) and (1.5,2) .. (2,2) .. controls (2.6,1.95) and (3.45,-0.5) .. (4,-0.5) .. controls (4.5,-0.5) and (5.4,1) .. (6,1) ;

\draw[->,Mapcolor] (0.2,0.1) -- (0.9,0.4);
\draw[->,Mapcolor] (2,2) -- (3.5,1.7);
\foreach \x/\g in {A/150,C/60}
\draw[fill=white](\x) circle(0.05)+(\g:0.3)node{$\x$};
\end{tikzpicture}\end{hinhanh}}
\loigiai{

}
\end{ex}
\begin{ex}%Tailieuchuan
\immini[thm]{Một chất điểm chuyển động có quỹ đạo là đường nét đứt như hình bên. Đâu là vectơ vận tốc trung bình của chất điểm khi chuyển động từ $A$ đến $B$?
\choice
{\True $\vec v_1$}
{$\vec v_2$}
{$\vec v_3$}
{$\vec v_A+\vec v_B$}}{\begin{tikzpicture}[scale=1,thick,>=stealth',draw=Mapcolor]
\path
(0,0)coordinate(O)
(0.3,0.2)coordinate(A) 
(4.9,0.2)coordinate(B) 
;
\draw[dashed,Mapcolor] (0,0) .. controls (0.7,0.2) and (1.5,2) .. (2,2) .. controls (2.6,1.95) and (3.45,-0.5) .. (4,-0.5) .. controls (4.5,-0.5) and (5.4,1) .. (6,1) .. controls (6.65,1) and (7,0.5) .. (7.5,0);

\draw[->,Mapcolor] (0.3,0.2) -- (1.1,0.8);
\draw[->,Mapcolor] (1.1,0.8) -- (6.4,1.8);
\draw[->,Mapcolor] (4.9,0.2) -- (6.4,1.8);
\foreach \x/\g in {A/150,B/90}
\draw[fill=white](\x) circle(0.05)+(\g:0.3)node{$\x$};
\draw[->,Mapcolor] (0.3,0.2) -- (1.4,0.2);
\draw[->,Mapcolor] (4.9,0.2) -- (1.4,0.2);
\node at (0.3,0.7) {$\vec v_A$};
\node at (5.9,0.7) {$\vec v_B$};
\node at (0.7,0) {$\vec v_1$};
\node at (2.8,0) {$\vec v_2$};
\node at (3.7,1.5) {$\vec v_3$};
\end{tikzpicture}}
\loigiai{

}
\end{ex}
\begin{ex}%Tailieuchuan
Khi vật chuyển động thẳng đều, vận tốc của vật có đặc điểm gì?
\choice
{\True Không thay đổi}
{Tăng dần theo thời gian}
{Giảm dần theo thời gian}
{Luôn bằng 0}
\loigiai{

}
\end{ex}
\begin{ex}%Tailieuchuan
Khi vật chuyển động thẳng đều theo chiều dương đã chọn, đồ thị độ dịch chuyển - thời gian có đặc điểm gì?
\choice
{Đường thẳng nằm ngang}
{Đường cong parabol}
{\True Đường thẳng xiên góc hướng lên}
{Đường thẳng xiên góc hướng xuống}
\loigiai{

}
\end{ex}
\begin{ex}%Tailieuchuan
Độ dốc của đồ thị độ dịch chuyển - thời gian trong chuyển động thẳng đều cho biết gì?
\choice
{\True Tốc độ của vật}
{Quãng đường đi được}
{Thời gian chuyển động}
{Hướng chuyển động}
\loigiai{

}
\end{ex}
\begin{ex}%Tailieuchuan
Khi độ dốc của đồ thị độ dịch chuyển - thời gian trong chuyển động thẳng đều bằng $0$, điều này có nghĩa là gì?
\choice
{Vật đang chuyển động nhanh dần}
{Vật đang chuyển động chậm dần}
{\True Vật đang đứng yên}
{Vật đang chuyển động ngược chiều dương}
\loigiai{}
\end{ex}
\begin{ex}%Tailieuchuan
Độ dịch chuyển của vật trong chuyển động thẳng đều là một hàm số của yếu tố nào?
\choice
{Tốc độ}
{\True Thời gian}
{Quỹ đạo}
{Tốc độ và Quỹ đạo}
\loigiai{

}
\end{ex}
\begin{ex}%book%[1P1K3-1][Loai1]
Lúc $7$ giờ, một người ở $A$ chuyển động thẳng đều với tốc độ không đổi $36\ km/h$ đuổi theo người ở $B$ đang chuyển động với tốc độ $5\ m/s$. Biết $AB = 18\ km$. Hai người gặp nhau cách $A$ bao xa?
\choice
{\True $36\ km$}
{$23\ km$}
{$31\ km$}
{$44\ km$}
\haidong{6}
\loigiai{
\begin{itemize}
\item Chọn gốc tọa độ tại A, gốc thời gian \bth{(t=0)} là lúc hai người bắt đầu xuất phát, chiều dương hướng từ A đến B.
\begin{itemize}
\item Phương trình chuyển động của người đi từ A: 
$x_A=x_{0A}+v_A.t=36t$.
\item Phương trình chuyển động của người đi từ B: 
$x_B=x_{0B}+v_{B}t=18+18t$.
\end{itemize}
\item Xác định vị trí và thời điểm hai xe gặp nhau.
\begin{itemize}
\item Lúc hai xe gặp nhau: \bth{x_A=x_B}
$\Leftrightarrow 36t=18+18t\Leftrightarrow 18t=18\Rightarrow t=1h$.
\item Thay \bth{t=1h} vào \bth{x_A} ta được \bth{x_A=36.1=36\ km}. 
\end{itemize}
\item Vậy hai người gặp nhau lúc 8h và cách gốc tọa độ A 36 km.
\end{itemize}
}
\haidong{7} \end{ex}
\begin{ex}%book%[1P1K3-1][Loai1]
Lúc $8$ giờ sáng, xe $1$ khởi hành từ $A$ chuyển động thẳng đều về $B$ với tốc độ $10\ m/s$. Nửa giờ sau, xe $2$ chuyển động thẳng đều từ $B$ đến $A$ và gặp nhau lúc $9$ giờ $30$ phút. Biết $AB = 72\ km$. Lúc hai xe cách nhau $13,5\ km$ sau khi xe $1$ vượt qua xe $2$ là mấy giờ?
\choice
{$9$ giờ $05$ phút}
{\True $9$ giờ $45$ phút}
{$9$ giờ $55$ phút}
{$10$ giờ $45$ phút}
\haidong{14}
\loigiai{
\begin{itemize}
\item Chọn gốc tọa độ tại A, gốc thời gian là lúc xe 1 bắt đầu xuất phát, chiều dương từ A đến B.
\item Phương trình chuyển động của xe 1:
\begin{align*}
x_{1}=x_{01}+v_1t=36t.
\end{align*}
\item Phương trình chuyển động của xe 2:
\begin{align*}
x_2=x_{02}+v_2(t-0,5)=72-v_2(t-0,5).
\end{align*}
\item Hai xe gặp nhau lúc 9 giờ 30 phút, tức là \bth{t=1,5h}, ta có
\begin{align*}
36t=72-v_2(t-0,5)\Leftrightarrow 36.2=72-v_2(1,5-0,5)\Rightarrow v_2=18\ km/h.
\end{align*}
\item Vậy phương trình chuyển động của xe 2: \bth{x_2=72-18(t-0,5)=81-18t}.
\item Khoảng cách giữa hai xe:
\begin{align*}
\abs{x_1-x_2}=13,5&\Leftrightarrow \abs{36t-(81-18t)}=\abs{54t-81}=13,5\\
&\Leftrightarrow\slbrace{\begin{array}{ll}
54t-81 & =13,5. \\ 
54t-81 & =-13,5.
\end{array}}\tag{1}
\end{align*}
\item Giải $(1)$ ta được \bth{t=1,25h} tức là 9 giờ 15 phút hoặc \bth{t=1,75h} tức là 9 giờ 45 phút.
\end{itemize}
}
\haidong{7} \end{ex}









\subsubsection*{PHẦN 2. Thí sinh trả lời từ câu 1 đến câu 4. Trong mỗi ý a), b), c), d) ở mỗi câu, thí sinh chọn đúng hoặc sai.}
\setcounter{ex}{0}

\begin{ex}%Tailieuchuan
Một chiếc ô tô xuất phát từ A lúc 6 giờ sáng, chuyển động thẳng đều theo hướng Bắc tới B, cách A $180\ km$. Xe tới B lúc 8 giờ 30 phút. Sau 30 phút đỗ tại B, xe chạy ngược về A với tốc độ $60 \ km / h$.
\choiceTF[t]
{Ô tô sẽ về tới A lúc 11h} 
{Tốc độ trung bình của ô tô trong cả quá trình là $70 \ km / h$} 
{Độ dịch chuyển của xe từ 6 h đến 8 h 30 là $160\ km$} 
{\True Độ dịch chuyển của xe đã đi được đến lúc 9 h 30 phút là $150\ km$} 
\loigiai{ 
\begin{itemchoice}
\itemch
\itemch
\itemch
\itemch
\end{itemchoice}
}
\end{ex}
\begin{ex}%Tailieuchuan
Lúc $6\ h$ sáng, xe máy đi từ $A$ đến $B$ theo hướng Tây - Đông với tốc độ $30 \ km/h$ và xe ô tô đi từ $A$ đến $C$ theo hướng Nam - Bắc với tốc độ $50 \ km/h$. Chọn gốc thời gian $t_0=0$ là lúc $6\ h$ sáng.
\choiceTF[t]
{\True Phương trình chuyển động của xe ô tô và xe máy theo thời gian là $d_x=30t$ và $d_y=50t$} 
{Độ dịch chuyển của ô tô so với xe máy là $d=30 t+50 t$} 
{Độ dịch chuyển của ô tô so với xe máy lúc $7$ giờ $30$ phút là $120\ km$} 
{\True Giả sử rằng vào lúc $7$ giờ $30$ phút ô tô và xe máy đi vào đường thẳng $BC$ nhưng ngược chiều nhau với vận tốc không đổi thì sau khoảng $1,093 \ h$ sẽ gặp nhau} 
\loigiai{ 
\begin{itemchoice}
\itemch
\itemch
\itemch
\itemch
\end{itemchoice}
}
\end{ex}
\begin{ex}%Tailieuchuan
Hai người ở hai đầu một đoạn đường thẳng $AB$ dài $10\ km$ đi bộ ngược chiều đến gặp nhau. Người ở $A$ đi trước người ở $B$ $0,5$ giờ. Sau khi người ở $B$ đi được $1$ giờ thì hai người gặp nhau. Biết hai người đi nhanh như nhau.
\choiceTF[t]
{Tốc độ của hai người ở $A$ và $B$ lần lượt là: $v_A=8\ km/h$ và $v_B=-8\ km/h$} 
{\True Vectơ vận tốc của hai người ngược chiều nhau} 
{\True Hai người gặp nhau sau khi người xuất phát từ $A$ đi được $1,5\ h$ tại vị trí cách $A$ một khoảng là $d_A = 4.1,5 = 6\ km$.} 
{Sau $2h$ kể từ khi $A$ bắt đầu đi khoảng cách của hai người là $8\ km$} 
\loigiai{ 
\begin{itemchoice}
\itemch
\itemch
\itemch
\itemch
\end{itemchoice}
}
\end{ex}
\begin{ex}
\immini[thm]{Một xe mô hình nhỏ chuyển động giữa hai vị trí $B$, $C$ trên một đoạn đường thẳng. Ban đầu, xe xuất phát từ $A$, đi về phía $B$. Đồ thị độ dịch chuyển - thời gian của xe được cho như hình vẽ, các giai đoạn $(1)$, $(5)$ có dạng parabol, giai đoạn $(2)$, $(3)$, $(4)$ có dạng đường thẳng.}{\begin{tikzpicture}[scale=1,thick,>=stealth',draw=Mapcolor]
\path
(0,0)coordinate(O)
(6,0)coordinate(x)
(0,4)coordinate(y)
(1,5)coordinate(a)
(1.5,5)coordinate(b)
(3.5,-0.5)coordinate(c)
(4,-0.5)coordinate(d)
(5,0)coordinate(e)

;
\draw  [dashed,help  lines, xstep=0.5cm,ystep=1cm]  (0,-1)  grid  (5,5);
\draw [->](O)--(x);

\draw[->] (0,-1) -- (0,5.5);
\draw[dashed,help  lines] (0,-0.5) -- (5,-0.5);

\foreach  \x  in  {2,3,6,7,8,10}
{\pgfmathtruncatemacro{\label}{10*\x};
\node[below]
(\label) at (.5*\x,0) {\label};
\draw[fill=white](.5*\x,0)circle(0.5pt);}

\foreach  \y  in  {-1,0,2,4,6,8,10}
{\pgfmathtruncatemacro{\label}{10*\y};
\node[left]
(\label) at (0,.5*\y) {\label};
\draw[fill=white](0,.5*\y)circle(0.5pt);}


\node at (0,6) {$d\ (m)$};
\node at (6.5,0) {$t\ (s)$};
\draw[Mapcolor,very thick] (0,0) .. controls (0.5,2) and (0.8,3.4) .. (1,5);
\draw[Mapcolor,very thick] (1,5) -- (1.5,5) -- (3.5,-0.5) -- (4,-0.5);
\draw[Mapcolor,very thick] (4,-0.5) .. controls (4.5,-0.5) and (4.9,-0.5) .. (5,0);
\foreach \x/\g in {a/90,b/180,c/0,d/0,e/0}
\draw[fill=white](\x) circle(0.05)+(\g:0.3)node{};
\node at (0.9,2.6) {(1)};
\node at (1.3,5.3) {(2)};
\node at (2.9,2.4) {(3)};
\node at (3.7,-0.8) {(4)};
\node at (4.9,-0.7) {(5)};

\path
(1,7)coordinate(A)
(5,7)coordinate(B)
(0,7)coordinate(C)
;
\draw[Mapcolor,very thick] (C)--(B);
\foreach \x/\g in {A/90,B/90,C/90}
\draw[fill=white](\x) circle(0.05)+(\g:0.3)node{$\x$};
\end{tikzpicture}}
\choiceTF[t]
{Trong giai đoạn $(1)$, xe chuyển động thẳng đều}
{\True Trong cả quá trình chuyển động, xe đổi chiều chuyển động hai lần}
{Độ dịch chuyển của xe trong toàn bộ quá trình chuyển động là $220\ m$}
{\True Tốc độ trung bình của xe trong toàn bộ quá trình chuyển động là $2,2 \ m/s$}
\loigiai{
\begin{itemchoice}
\itemch 		\itemch 		
\itemch 		\itemch 		
\end{itemchoice}
}
\end{ex}

\subsubsection*{PHẦN 3. Thí sinh trả lời từ câu 1 đến câu 6.}
\setcounter{ex}{0}
\begin{ex}%Tailieuchuan
\immini[thm]{Hình vẽ là đồ thị độ dịch chuyển - thời gian của một chiếc ô tô chạy từ $A$ đến $B$ trên một đường thẳng. Tốc độ của xe là bao nhiêu $km/h$?
}{\begin{tikzpicture}[scale=1,thick,>=stealth',draw=Mapcolor]
\path
(0,0)coordinate(O)
(3.5,0)coordinate(x)
(0,3.5)coordinate(y)
(2.5,1.5)coordinate(A)
;
\draw  [dashed,help  lines, xstep=0.5cm,ystep=0.5cm]  (0,0)  grid  (3,3);
\draw [->](O)--(x);
\draw [->](O)--(y);

\node at (0,4) {$d\ (km)$};
\node at (4,0) {$t\ (h)$};
\node[left] at (0,1) {$60$};
\node[below] at (0.5,0) {$1$};
\draw[very thick,Mapcolor] (0,0) -- (3.3,2);
\foreach \x/\g in {A/90}
\draw[fill=white](\x) circle(0.05)+(\g:0.3)node{};
\end{tikzpicture}}
\shortans[oly]{18}
\loigiai{
}
\end{ex}
\begin{ex}%Tailieuchuan
\immini[thm]{Một chất điểm chuyển động trên một đường thẳng. Đồ thị độ dịch chuyển theo thời gian của chất điểm được mô tả trên hình vẽ. Tốc độ trung bình của chất điểm trong khoảng thời gian từ $0,5\ s$ đến $5\ s$ là bao nhiêu $cm/s$ (làm tròn kết quả đến chữ số hàng phần mười)?
}{\begin{tikzpicture}[scale=1,thick,>=stealth',draw=Mapcolor]
\path
(0,0)coordinate(O)
(0,-1.5)coordinate(O')
(6.5,0)coordinate(x)
(0,3)coordinate(y)
;
\draw  [dashed,help  lines, xstep=0.5cm,ystep=0.5cm]  (0,-1)  grid  (6,2.5);
\draw [->](O)--(x);
\draw [->](O')--(y);

\node at (0,3.5) {$d\ (cm)$};
\node at (7,0) {$t\ (s)$};

\foreach  \x  in  {1,2,...,4}
{\pgfmathtruncatemacro{\label}{5*\x};
\draw[fill=white](.5*\x,0)circle(0.5pt);}


\foreach  \y  in  {-2,-1,1,2,...,4}
{\pgfmathtruncatemacro{\label}{5*\y};
\draw[fill=white](0,.5*\y)circle(0.5pt);}
\node[below] at (1,0) {$1$};
\node[left] at (0,2) {$4$};
\draw[very thick,Mapcolor] (0,0) -- (1,2) -- (2.5,-1) -- (4,-1) -- (5,0);
\draw (0.5,-2);
\end{tikzpicture}}
\shortans[oly]{2,2}
\loigiai{
}
\end{ex}

\begin{ex}%Tailieuchuan
\immini[thm]{Đồ thị độ dịch chuyển - thời gian của hai chiếc xe $(1)$ và $(2)$ được biểu diễn như hình vẽ bên. Hai xe gặp nhau tại vị trí cách vị trí xuất phát của xe $(1)$ một đoạn bao nhiêu $km$?
}{\begin{hinhanh}{6}\begin{tikzpicture}[scale=1,thick,>=stealth',draw=Mapcolor]
\path
(0,0)coordinate(O)
(3,0)coordinate(x)
(0,3)coordinate(y)
(1,1)coordinate(A)
;
\draw  [dashed,help  lines, xstep=0.5cm,ystep=0.5cm]  (0,0)  grid  (2.5,2.5);
\draw [->](O)--(x);
\draw [->](O)--(y);

\node at (0,3.5) {$d\ (km)$};
\node at (3.5,0) {$t\ (h)$};
\foreach  \x  in  {1,2,...,4}
{\pgfmathtruncatemacro{\label}{5*\x};
\draw[fill=white](.5*\x,0)circle(0.5pt);}


\foreach  \y  in  {1,2,...,4}
{\pgfmathtruncatemacro{\label}{5*\y};
\draw[fill=white](0,.5*\y)circle(0.5pt);}
\draw[dashed,thin] (0,1) -- (1,1) -- (1,0);
\node[left] at (0,2) {$80$};
\draw[Mapcolor] (0,2) -- (2,0);
\draw[Mapcolor] (0,0) -- (2.5,2.5);
\node at (0.7,1.7) {(1)};
\node at (1.6,2) {(2)};
\node[below] at (1,0) {$2$};
\foreach \x/\g in {A/90}
\draw[fill=white](\x) circle(0.05)+(\g:0.3)node{};
\end{tikzpicture}\end{hinhanh}}
\shortans[oly]{40}
\loigiai{
}
\end{ex}
\begin{ex}%Tailieuchuan
\immini[thm]{Hai chiếc xe $A$ $(I)$ và $B$  $(II)$ chuyển động trên một đoạn đường thẳng, với độ dịch chuyển tại các thời điểm khác nhau được được mô tả như hình vẽ. Tỉ số quãng đường của xe A so với xe B là bao nhiêu?}{\begin{hinhanh}{6}
\begin{tikzpicture}[scale=1,thick,>=stealth',draw=Mapcolor]
\path
(0,0)coordinate(O)
(0,-1.5)coordinate(O')
(3.5,0)coordinate(x)
(0,3)coordinate(y)
;
\draw  [dashed,help  lines, xstep=0.5cm,ystep=0.5cm]  (0,-1)  grid  (3,2.5);
\draw [->](O)--(x);
\draw [->](O')--(y);

\node at (0,3.5) {$d\ (km)$};
\node at (4,0) {$t\ (h)$};

\foreach  \x  in  {1,2,...,4}
{\pgfmathtruncatemacro{\label}{5*\x};
\draw[fill=white](.5*\x,0)circle(0.5pt);}


\foreach  \y  in  {-2,-1,1,2,...,4}
{\pgfmathtruncatemacro{\label}{5*\y};
\draw[fill=white](0,.5*\y)circle(0.5pt);}
\node[below] at (1,0) {$1$};
\node[below] at (2,0) {$2$};
\node[left] at (0,1) {$60$};
\node[left] at (0,0.5) {$30$};
\node[left] at (0,0) {$O$};
\node[left] at (0,-1) {$-60$};
\draw[very thick,Mapcolor] (0,0) -- (1,-1) -- (3.5,1.5);

\draw[very thick,Mapcolor] (0,0) -- (1,0.5) -- (2,0.5) -- (4,1.5);
\node at (2.9,0.5) {(I)};
\node at (4.1,1.1) {(II)};
\end{tikzpicture}\end{hinhanh}}
\shortans[oly]{2}
\loigiai{
}
\end{ex}


\begin{ex}%Tailieuchuan
Lúc $7h$, ô tô thứ nhất đi qua điểm $A$, ô tô thứ hai đi qua điểm $B$ cách $A$ $10\ km$. Xe đi qua $A$ với vận tốc $50\ km/h$, xe đi qua $B$ với vận tốc $40\ km/h$. Biết hai xe chuyển động cùng chiều theo hướng từ $A$ đến $B$. Coi chuyển động của hai ô tô là chuyển động đều. Thời điểm hai xe gặp nhau lúc mấy giờ?
\shortans[oly]{8}
\loigiai{
}
\end{ex}
\begin{ex}%Tailieuchuan
Hai xe coi là chuyển động thẳng đều từ $A$ đến $B$ cách nhau $60\ km$. Xe $(1)$ có tốc độ $15\ km/h$ và chạy liên tục không nghỉ. Xe $(2)$ khởi hành sớm hơn $1$ giờ nhưng dọc đường phải dừng lại $2$ giờ. Xe $(2)$ phải có tốc độ bao nhiêu $km/h$ để tới $B$ cùng lúc với xe $(1)$?
\shortans[oly]{20}
\loigiai{}
\end{ex}







\newpage 
\TieudeT{PHONG TỎA VẬT LÍ}
\subsubsection*{PHẦN I. Thí sinh trả lời câu hỏi từ câu 1 đến câu 18. Mỗi câu hỏi thí sinh chỉ chọn một phương án}
\setcounter{ex}{0}

\begin{ex}%book
Đồ thị độ dịch chuyển theo thời gian của chuyển động thẳng đều là một
\choice
{đường parabol}
{đường hypebol}
{\True đoạn thẳng}
{hình tròn}

\loigiai{
}
\haidong{7} \end{ex}
\begin{ex}%book
Trong đồ thị độ dịch chuyển $(d)$ - thời gian $(t)$, đường thẳng song song với trục thời gian $(Ot)$ mô tả một vật
\choice
{\True đang đứng yên}
{chuyển động thẳng nhanh dần đều}
{chuyển động thẳng chậm dần đều}
{chuyển động thẳng đều}

\loigiai{
}
\haidong{7} %<MyLT1>
\end{ex}
\begin{ex}%book
\immini[thm]{Cho đồ thị độ dịch chuyển - thời gian của một vật như hình. Phát biểu nào sau đây đúng?
\choice
{Vật chuyển động thẳng đều theo chiều dương rồi dừng lại}
{Vật chuyển động thẳng đều theo chiều âm sau đó đổi chiều chuyển động ngược lại}
{Vật lúc đầu đứng yên sau đó chuyển động theo chiều dương}
{\True Vật chuyển động thẳng đều theo chiều dương rồi đổi chiều chuyển động ngược lại}}{\begin{tikzpicture}[very thick,draw=Mapcolor]
\tkzDefPoints{0/0/o, 0/2.5/d, 4/0/t}
\tkzDrawSegments[->](o,d o,t)
\draw[Mapcolor, very thick] (0,0) -- (2,2) -- (3,1);
\draw [dashed](2,2) -- (2,0);
\tkzLabelPoint[below left](o){o}
\tkzLabelPoint[above right](d){d}
\tkzLabelPoint[below](t){t}
\end{tikzpicture}}

\loigiai{
}
\haidong{7} \end{ex}
\begin{ex}%book%7.3
Cặp đồ thị nào dưới đây là của chuyển động thẳng đều?
\begin{center}
\begin{tikzpicture}[scale=1,thick,>=stealth',draw=Mapcolor]
\tkzDefPoints{0/0/o, 0/3/y, 3/0/x}
\tkzDefShiftPoint[o](45:3cm){a}
\tkzDrawSegments[->](o,x o,y)
\tkzDrawSegments[very thick,Mapcolor](o,a)
\tkzLabelPoint[below left](o){$O$}
\tkzLabelPoint[below](x){$t$}
\tkzLabelPoint[left](y){$d$}
\node at (1.5,-0.5){(I)};
\begin{scope}[shift={(4,0)},rotate=0]
\tkzDefPoints{0/0/o, 0/3/y, 3/0/x}
\tkzDefShiftPoint[o](45:3cm){a}
\tkzDrawSegments[->](o,x o,y)
\tkzLabelPoint[below left](o){$O$}
\tkzLabelPoint[below](x){$t$}
\tkzLabelPoint[left](y){$d$}
\draw[very thick,Mapcolor](0,2)--(3,2);
\node at (1.5,-0.5){(II)};
\end{scope}
\begin{scope}[shift={(8,0)},rotate=0]
\tkzDefPoints{0/0/o, 0/3/y, 3/0/x}
\tkzDefShiftPoint[o](45:3cm){a}
\tkzDrawSegments[->](o,x o,y)
\tkzDrawSegments[very thick,Mapcolor](o,a)
\tkzLabelPoint[below left](o){$O$}
\tkzLabelPoint[below](x){$t$}
\tkzLabelPoint[left](y){$v$}
\node at (1.5,-0.5){(III)};
\end{scope}
\begin{scope}[shift={(12,0)},rotate=0]
\tkzDefPoints{0/0/o, 0/3/y, 3/0/x}
\tkzDefShiftPoint[o](45:3cm){a}
\tkzDrawSegments[->](o,x o,y)
\tkzLabelPoint[below left](o){$O$}
\tkzLabelPoint[below](x){$t$}
\tkzLabelPoint[left](y){$v$}
\draw[very thick,Mapcolor](0,2)--(3,2);
\node at (1.5,-0.5){(IV)};
\end{scope}
\end{tikzpicture}
\end{center}
\choice
{I và III}
{\True I và IV}
{II và III}
{II và IV}
\haidong{7} \end{ex}
\begin{ex}%book %[1P1B4-1][Loai1] 
Trong các đồ thị sau, đồ thị nào mô tả chuyển động thẳng đều?
\begin{center}
\begin{tikzpicture}[scale=1,thick,>=stealth',x=1.5cm,y=0.2cm,draw=Mapcolor]
\tkzDefPoints{0/0/o, 0/12/x, 2/0/t, 0/4/m, 2/10/n, 2/0/n1, 0/12/n2}
\tkzDrawLines[very thick,color=Mapcolor,add = 0 and 0](m,n)
\tkzDrawSegments[->](o,x o,t)
\tkzLabelPoint[below](o){$O$}
\tkzLabelPoint[right](x){$x$}
\tkzLabelPoint[above](t){$t$}
\node at (1,-1.5){(a)};
\begin{scope}[shift={(3,0)}]
\tkzDefPoints{0/5/o', 0/0/o, 0/12/x, 2/5/t, 0/5/m, 2/0/n, 2/0/n1, 0/12/n2}
\tkzDrawLines[very thick,color=Mapcolor,add = 0 and 0](m,n)
\tkzDrawSegments[->](o,x o',t)
\tkzLabelPoint[left](o'){$O$}
\tkzLabelPoint[right](x){$x$}
\tkzLabelPoint[above](t){$t$}
\node at (1,-1.5){(b)};
\end{scope}
\begin{scope}[shift={(6,0)}]
\tkzDefPoints{0/0/o, 0/12/x, 2/0/t, 0/5/m, 2/5/n, 2/0/n1, 0/12/n2}
\tkzDrawLines[very thick,color=Mapcolor,add = 0 and 0](m,n)
\tkzDrawSegments[->](o,x o,t)
\tkzLabelPoint[below](o){$O$}
\tkzLabelPoint[right](x){$v$}
\tkzLabelPoint[above](t){$t$}
\node at (1,-1.5){(c)};
\end{scope}
\begin{scope}[shift={(9,0)}]
\tkzDefPoints{0/5/o', 0/0/o, 0/12/x, 2/5/t, 0/2/m, 2/2/n, 2/0/n1, 0/12/n2}
\tkzDrawLines[very thick,color=Mapcolor,add = 0 and 0](m,n)
\tkzDrawSegments[->](o,x o',t)
\tkzLabelPoint[left](o'){$O$}
\tkzLabelPoint[right](x){$v$}
\tkzLabelPoint[above](t){$t$}
\node at (1,-1.5){(d)};
\end{scope}
\end{tikzpicture}
\end{center}
\choice
{Đồ thị $(a)$}
{Đồ thị $(b)$ và $(d)$}
{Đồ thị $(a)$ và $(c)$}
{\True Các đồ thị $(a)$, $(b)$, $(c)$ và $(d)$ đều đúng}
\loigiai{
\begin{itemize}
\item Đồ thị (a): do x(t) là một đường thẳng $\Rightarrow$mô tả chuyển động thẳng đều.
\item Đồ thị (b): Tương tự như (a)
\item Đồ thị (c): Do $v = const \Rightarrow$ mô tả chuyển động thẳng đều.
\item Đồ thị (d): Tương tự như (c)
\end{itemize}	}
\haidong{7} \end{ex}
\begin{ex}%book
\immini[thm]{Hình bên là đồ thị độ dịch chuyển của vật chuyển động thẳng đều. Dựa vào đồ thị cho biết độ dịch chuyển của vật từ thời điểm $0\ s$ đến $3\ s$ là
\choice
{$80\ m$}
{$160\ m$}
{\True $240\ m$}
{$320\ m$}}{\begin{tikzpicture}[very thick,draw=Mapcolor]
\tkzDefPoints{0/0/o, 0/3/d, 4/0/t}
\tkzDrawSegments[->, thick](o,d o,t)
\draw [dashed, thin] (0,0.6) -- (0.8,0.6) -- (0.8,0);
\draw[dashed, thin] (0,1.2) -- (1.6,1.2) -- (1.6,0);
\draw[dashed, thin] (0,1.8) -- (2.4,1.8) -- (2.4,0);
\draw [dashed, thin](0,2.4) -- (3.2,2.4)  -- (3.2,0);
\draw [Mapcolor, very thick](0,0) -- (3.2,2.4);
\tkzLabelPoint[below](o){$O$}
\tkzLabelPoint[right](d){$d~(m)$}
\tkzLabelPoint[below](t){$t~(s)$}
\node [left]at (0,0.6) {80};
\node [below]at (1.6,0) {2};
\node[below] at (3.2,0) {4};
\end{tikzpicture}
}
\loigiai{

}
\haidong{7} \end{ex}
\begin{ex}%book
\immini[thm]{Hình vẽ bên là đồ thị độ dịch chuyển - thời gian của một chiếc xe chạy từ $A$ đến $B$ trên đường một đường thẳng. Xe này có tốc độ là
\choice
{\True $30\ km/h$}
{$60\ km/h$}
{$15\ km/h$}
{$45\ km/h$}}{\begin{hinhanh}{6}\begin{tikzpicture}[very thick, scale=1.2, draw=Mapcolor]
\tkzDefPoints{0/0/o, 0/3/d, 3/0/t}
\tkzDrawSegments[->, thick](o,d o,t)
\draw  [help  lines, dashed, xstep=0.5cm,ystep=0.5cm]  (0,0)  grid  (2.5,2.5);
\draw[very thick,Mapcolor] (0.5,0.5) node [below left, scale=0.9] {A} -- (2.5,2.5) [above, scale= 0.9]node {B};
\node [left, scale=0.8]at (0,0.5) {30};
\node [left, scale=0.8]at (0,1) {60};
\node[left, scale=0.8] at (0,1.5) {90};
\node[left, scale=0.8] at (0,2) {120};
\node[left, scale=0.8] at (0,2.5) {150};
\node [below, scale=0.8]at (0.5,0) {1};
\node [below, scale=0.8]at (1,0) {2};
\node [below, scale=0.8]at (1.5,0) {3};
\node [below, scale=0.8]at (2,0) {4};
\node [below, scale=0.8]at (2.5,0) {5};
\tkzLabelPoint[below left](o){O}
\tkzLabelPoint[above](d){$d~(km)$}
\tkzLabelPoint[above right](t){$t~(h)$}
\end{tikzpicture}\end{hinhanh}
}
\haidong{1}
\loigiai{
}
\haidong{7} \end{ex}
\begin{ex}%book
\immini[thm]{Hình bên là đồ thị độ dịch chuyển - thời gian của một vật chuyển động. Tại thời điểm $t=4\ s$, vật cách vị trí ban đầu (vị trí ở $t=0$) một khoảng bằng
\choice
{$40\ m$}
{\True $20\ m$}
{$50\ m$}
{$10\ m$}}{\begin{hinhanh}{6}\begin{tikzpicture}[very thick,draw=Mapcolor]
\tkzDefPoints{0/0/o, 0/2.5/d, 3/0/t}
\tkzDrawSegments[->, thick](o,d o,t)
\draw  [help  lines, black, dashed, xstep=0.5cm,ystep=0.5cm]  (0,0) node (v1) {}  grid  (2.5,2);
\draw [Mapcolor, very thick](0,0) -- (1,1.5) -- (2,1);
\tkzLabelPoint[below left](o){O}
\tkzLabelPoint[above](d){$d~(m)$}
\tkzLabelPoint[above right](t){$t~(s)$}

\node [left] at (0,0.5) {10};
\node [left]at (0,1.5) {30};
\node [below]at (1,0) {2};
\node[below] at (2,0) {4};
\end{tikzpicture}\end{hinhanh}}
\loigiai{
}
\haidong{7} \end{ex}
\begin{ex}%Tailieuchuan
Nếu đồ thị độ dịch chuyển - thời gian là một đường thẳng nằm ngang, điều này có nghĩa là gì?
\choice
{Vật đang chuyển động đều theo chiều dương}
{Vật đang chuyển động đều theo chiều âm}
{\True Vật đang đứng yên}
{Vật đang chuyển động với vận tốc không đổi}
\loigiai{

}
\end{ex}
\begin{ex}%Tailieuchuan
Khi độ dịch chuyển của vật tăng đều theo thời gian, điều này có nghĩa là gì về độ dốc của đồ thị độ dịch chuyển - thời gian?
\choice
{Độ dốc giảm dần}
{\True Độ dốc không đổi}
{Độ dốc tăng dần}
{Độ dốc bằng $0$}
\loigiai{

}
\end{ex}
\begin{ex}%Tailieuchuan
Để xác định vận tốc của vật từ đồ thị độ dịch chuyển - thời gian trong chuyển động thẳng đều, ta cần làm gì?
\choice
{Tính diện tích dưới đồ thị}
{\True Tính độ dốc của đồ thị}
{Tính khoảng cách giữa hai điểm trên đồ thị}
{Tính giá trị lớn nhất của đồ thị}
\loigiai{

}
\end{ex}
\begin{ex}%Tailieuchuan
Nếu một vật chuyển động thẳng đều và đổi chiều tại một điểm, đồ thị độ dịch chuyển - thời gian sẽ có đặc điểm gì tại điểm đó?
\choice
{Đồ thị sẽ có dạng parabol}
{Đồ thị sẽ có một đoạn thẳng nằm ngang}
{Đồ thị sẽ có một đoạn thẳng thẳng đứng}
{\True Đồ thị sẽ có một điểm gấp khúc}
\loigiai{

}
\end{ex}
\begin{ex}%Tailieuchuan
Khi đồ thị độ dịch chuyển - thời gian có độ dốc lớn, điều này có nghĩa là gì về tốc độ của vật?
\choice
{Tốc độ của vật nhỏ}
{\True Tốc độ của vật lớn}
{Vật đang đưng yên}
{Vật đang tăng tốc}

\loigiai{

}
\end{ex}
\begin{ex}
Bố bạn A đưa A đi học bằng xe máy vào lúc 7 giờ. Sau 5 phút xe đạt tốc độ 30 km/h. Sau 10 phút nữa, xe tăng tốc lên thêm 15 km/h. Đến gần trường, xe giảm dần tốc độ và dừng trước cổng trường lúc 7 giờ 30 phút. Biết quãng đường từ nhà đến trường dài 15 km. Tốc độ trung bình của xe máy chở A khi đi từ nhà đến trường là

\choice
{\True $30\ km/$giờ}
{$150\ km/$giờ}
{$120\ km /$giờ}
{$100\ km/$giờ}
\loigiai{

}
\end{ex}
\begin{ex}%Tailieuchuan
\immini[thm]{Hình dưới là đồ thị độ dịch chuyển - thời gian của hai vật chuyển động thẳng cùng hướng. Tỉ lệ vận tốc $v_{A}$: $v_{B}$ là
\choice[2]
{$3:1$}
{\True $1:3$}
{$\sqrt{3}:1$}
{$1:\sqrt{3}$}}{\begin{hinhanh}{6}\begin{tikzpicture}[line cap=round, line join=round, scale=1, font=\footnotesize, >=stealth,draw=Mapcolor]
\path
(0,0)coordinate(O)
(3,0)coordinate(x)
(0,3)coordinate(y)
(2.5,1.5)coordinate(A)
(1.5,2)coordinate(B)
;
\draw [->](O)--(x);
\draw [->](O)--(y);

\node at (0,3.5) {$d$};
\node at (3.5,0) {$t$};

\draw[very thick,Mapcolor] (O)--(A);
\draw[very thick,Mapcolor] (O)--(B);

\draw pic[draw,angle radius=0.6cm]{angle=x--O--A};
\draw pic[draw,angle radius=0.55cm]{angle=x--O--A};

\draw pic[draw,angle radius=0.4cm]{angle=x--O--B};

\foreach \x/\g in {A/90,B/90}
\draw[fill=white](\x) circle(0.0)+(\g:0.3)node{$\x$};
\node at (0.9,0.7) {$60^\circ$};
\node at (1,0.3) {$30^\circ$};
\end{tikzpicture}\end{hinhanh}
}
\loigiai{

}
\end{ex}
\begin{ex}%Tailieuchuan
Cặp đồ thị nào ở hình dưới đây là của chuyển động thẳng đều?
\begin{center}
\begin{tikzpicture}[scale=1,thick,>=stealth',draw=Mapcolor]
\tkzDefPoints{0/0/o, 0/3/y, 3/0/x}
\tkzDefShiftPoint[o](45:3cm){a}
\tkzDrawSegments[->](o,x o,y)
\tkzDrawSegments[very thick,Mapcolor](o,a)
\tkzLabelPoint[below left](o){$O$}
\tkzLabelPoint[below](x){$t$}
\tkzLabelPoint[left](y){$d$}
\node at (1.5,-0.5){(I)};
\begin{scope}[shift={(4,0)},rotate=0]
\tkzDefPoints{0/0/o, 0/3/y, 3/0/x}
\tkzDefShiftPoint[o](45:3cm){a}
\tkzDrawSegments[->](o,x o,y)
\tkzLabelPoint[below left](o){$O$}
\tkzLabelPoint[below](x){$t$}
\tkzLabelPoint[left](y){$d$}
\draw[very thick,Mapcolor](0,2)--(3,2);
\node at (1.5,-0.5){(II)};
\end{scope}
\begin{scope}[shift={(8,0)},rotate=0]
\tkzDefPoints{0/0/o, 0/3/y, 3/0/x}
\tkzDefShiftPoint[o](45:3cm){a}
\tkzDrawSegments[->](o,x o,y)
\tkzDrawSegments[very thick,Mapcolor](o,a)
\tkzLabelPoint[below left](o){$O$}
\tkzLabelPoint[below](x){$t$}
\tkzLabelPoint[left](y){$v$}
\node at (1.5,-0.5){(III)};
\end{scope}
\begin{scope}[shift={(12,0)},rotate=0]
\tkzDefPoints{0/0/o, 0/3/y, 3/0/x}
\tkzDefShiftPoint[o](45:3cm){a}
\tkzDrawSegments[->](o,x o,y)
\tkzLabelPoint[below left](o){$O$}
\tkzLabelPoint[below](x){$t$}
\tkzLabelPoint[left](y){$v$}
\draw[very thick,Mapcolor](0,2)--(3,2);
\node at (1.5,-0.5){(IV)};
\end{scope}
\end{tikzpicture}
\end{center}
\choice
{$I$ và $III$}
{\True $I$ và $IV$}
{$II$ và $III$}
{$II$ và $IV$}
\loigiai{

}
\end{ex}
\begin{ex}%Tailieuchuan
Ta có $A$ cách $B$ $72\ km$. Lúc $7$ giờ $30$ phút, xe ô tô $(1)$ khởi hành từ $A$ chuyển động thẳng đều về $B$ với tốc độ $36\ km/h$. Nửa giờ sau, xe ô tô $(2)$ chuyển động thẳng đều từ $B$ đến $A$ và gặp nhau lúc $8$ giờ $30$ phút. Tốc độ của xe ô tô thứ hai là
\choice
{$70 \ km/h$}
{\True $72 \ km/h$}
{$73 \ km/h$}
{$74 \ km/h$}
\loigiai{

}
\end{ex}







\begin{ex}%book%[1P1K3-1][Loai1]  
Lúc $7$ giờ một ô tô xuất phát đi từ $A$ về $B$ với tốc độ $\dfrac{50}{3}\ m/s$ và cùng lúc đó một ô tô khác xuất phát từ $B$ về $A$ với tốc độ $50\ km/h$. Biết $A$ và $B$ cách nhau $220\ km$. Lấy $AB$ làm trục tọa độ, $A$ là gốc tọa độ, chiều dương từ $A$ đến $B$ và gốc thời gian là lúc $7$ giờ. Thời điểm hai xe gặp nhau là
\choice
{\True $9h$}
{$10h$}
{$11h$}
{$12h$}
\haidong{8}
\loigiai{
\begin{center}
\begin{tikzpicture}[scale=1,thick,>=stealth', draw=Mapcolor]
\tkzDefPoints{0/0/o, 5/0/b, 6/0/x}
\tkzDrawPoints[size=2,fill=black](o,b)
\tkzDrawSegments[->](o,x)
\draw[very thick,Mapcolor,->](o)--++(0:1.6)node[below,black]{$\vec{v}_A$};
\draw[very thick,Mapcolor,->](b)--++(180:1.8)node[below,black]{$\vec{v}_B$};
\tkzLabelPoint[above](o){$O$}
\tkzLabelPoint[below](o){$A$}
\tkzLabelPoint[below](b){$B$}
\tkzLabelPoint[below](x){$x$}
\end{tikzpicture}
\end{center}
\begin{itemize}
\item Theo cách chọn của đề ta có: $\begin{cases}
x_{0A}=0;\ v_A=\dfrac{50}3=60\ km/h.\\
x_{0B}=220;\ v_B=-50\ km/h.
\end{cases}$ 
\item Do đó phương trình chuyển động của hai xe xe là: $\begin{cases}
x_A=60t.\\
x_B=220-50t.
\end{cases}$ (x đo bằng km và t đo bằng h).
\item Khi hai xe gặp nhau: $x_A=x_B\Leftrightarrow 60t=220-50t\Rightarrow t=2h\Rightarrow x_A=120\ km$.
\item Vậy hai xe gặp nhau lúc 9 giờ.
\end{itemize}	
}
\haidong{7} \end{ex}
\subsubsection*{PHẦN 2. Thí sinh trả lời từ câu 1 đến câu 4. Trong mỗi ý a), b), c), d) ở mỗi câu, thí sinh chọn đúng hoặc sai.}
\setcounter{ex}{0}
\begin{ex}
\immini[thm]{
Đồ thị vận tốc - thời gian của một vật như trên hình vẽ.}{\begin{tikzpicture}[line cap=round, line join=round, scale=1, font=\footnotesize, >=stealth, draw=Mapcolor]
\path
(0,0)coordinate(O)
(0,-0.5)coordinate(O')
(5,0)coordinate(x)
(0,3)coordinate(y)

;
\draw [->](O)--(x);
\draw [->](O')--(y);
\node[left] at (O) {$0$};

\node at (0,3.5) {$v\ (m/s)$};
\node at (5.5,0) {$t\ (s)$};


\draw[dashed] (0,1.5) -- (1,1.5) -- (1,0);
\draw[dashed] (2.5,0) -- (2.5,1.5);
\draw[very thick,Mapcolor] (0,0) -- (1,1.5) -- (2.5,1.5) -- (4,0);
\node[left] at (0,1.5) {$10$};
\node[below] at (1,0) {$5$};
\node[below] at (2.5,0) {$12$};
\node[below] at (4,0) {$22$};
\end{tikzpicture}
}
\choiceTF[t]
{\True Từ thời điểm $0$ giây đến $5$ giây, vật chuyển động nhanh dần đều theo chiều dương}
{Từ thời điểm $5$ giây đến $12$ giây, vật đứng yên}
{Từ thời điểm $12$ giây đến $22$ giây, vật chuyển động chậm dần đều theo chiều dương với gia tốc là $1\ m/s$}
{\True Từ thời điểm $0$ giây đến $5$ giây, quãng đường vật đi được là $25\ m$}
\loigiai{
\begin{itemchoice}
\itemch
\itemch
\itemch
\itemch
\end{itemchoice}}
\end{ex}
\begin{ex}
\immini[thm]{Xe buýt chuyển động giữa các trạm xe ($A$, $B$, $C$), trên đường thẳng, biết rằng khi đến trạm dừng, xe buýt dừng lại đón khánh lên/xuống trong $3$ phút. Bỏ qua khoảng thời gian xe tăng, giảm tốc độ ở lân cận các trạm xe. Đồ thị độ dịch chuyển - thời gian của xe như hình dưới. Chọn mốc thời gian lúc xe bắt đầu xuất phát tại $A$.}{\begin{tikzpicture}[line cap=round, line join=round, scale=1, font=\footnotesize, >=stealth,draw=Mapcolor]
\path
(0,0)coordinate(O)
(0,-0.5)coordinate(O')
(8.4,0)coordinate(x)
(0,5.2)coordinate(y)
;
\path
(0,7)coordinate(A)
(1.5,7)coordinate(B)
(3,7)coordinate(C)
(4.5,7)coordinate(D)
;
\draw [->](O)--(x);
\draw [->](O)--(y);
\node[left] at (O) {$0$};

\node at (0,5.8) {$d\ (km)$};
\node at (9,0) {$t$ (phút)};


\foreach  \x  in  {15,33,51,69,87}
{\pgfmathtruncatemacro{\label}{\x};
\node[below left,scale=0.75]
(\label) at (.075*\x,0) {\label};
\draw[fill=white](.075*\x,0)circle(0.5pt);}

\foreach  \x  in  {18,36,54,72,90,105}
{\pgfmathtruncatemacro{\label}{\x};
\node[below ,scale=0.75]
(\label) at (.075*\x,0) {\label};
\draw[fill=white](.075*\x,0)circle(0.5pt);}

\foreach  \y  in  {1,2,3}
{\pgfmathtruncatemacro{\label}{10*\y};
\node[left]
(\label) at (0,1.5*\y) {\label};
\draw[fill=white](0,1.5*\y)circle(0.5pt);}
\draw[dashed] (0,1.5) -- (6.77,1.5);

\draw[dashed] (0,3) -- (5.4,3);

\draw[dashed] (0,4.5) -- (4.04,4.5);

\draw[dashed] (1.13,0) -- (1.14,1.5);

\draw[dashed] (1.36,0) -- (1.36,1.5);

\draw[dashed] (2.48,0) -- (2.48,3);

\draw[dashed] (2.7,0) -- (2.7,3);

\draw[dashed] (4.04,0) -- (4.04,4.5);

\draw[dashed] (3.82,0) -- (3.82,4.5);

\draw[dashed] (5.18,0) -- (5.18,3);

\draw[dashed] (5.4,0) -- (5.4,3);

\draw[dashed] (6.54,0) -- (6.54,1.5);

\draw[dashed] (6.76,0.02) -- (6.76,1.52);
\draw[very thick,Mapcolor] (O) -- (1.14,1.5) -- (1.37,1.5) -- (2.48,3) -- (2.71,3) -- (3.81,4.5) -- (4.03,4.5) -- (5.19,3) -- (5.4,3) -- (6.53,1.5) -- (6.77,1.5) -- (7.87,0);
\node at (-0.3,0.3) {(A)};
\node at (8,0.4) {(A)};
\node at (1.2,1.8) {(B)};
\node at (6.6,1.8) {(B)};
\node at (2.6,3.3) {(C)};
\node at (5.3,3.3) {(C)};
\node at (3.9,4.7) {(D)};

\draw[very thick,Mapcolor] (A)--(D);
\foreach \x/\g in {A/90,B/90,C/90, D/90}
\draw[fill=white](\x) circle(0.05)+(\g:0.3)node{$\x$};
\draw[->,red,very thick] (0,6.5) -- (1,6.5);
\draw[->,red,very thick]  (3.75,6.75) .. controls (4.75,6.75) and (4.75,6.25) .. (3.75,6.25);
\end{tikzpicture}
}
\choiceTF[t]
{Sau $1$ giờ $30$ phút, xe buýt trở lại trạm $A$}
{\True Tốc độ trung bình của xe buýt trong chuyển động giữa các trạm là $40 \ km/h$}
{Tốc độ trung bình của xe buýt trên cả hành trình từ $A$ đến $D$ xấp xỉ $43,3 \ km / h$}
{\True Trong lúc chuyển động từ trạm $D$ về trạm $A$, xe gặp sự cố ở trạm dừng $B$ nên phải dừng lại tại trạm $B$ trong thời gian $5$ phút (vượt quá $2$ phút so với yêu cầu của hành trình). Để xe về trạm $A$ đúng giờ thì tốc độ trung bình của xe trong giai đoạn chuyển động từ $B$ về $A$ xấp xỉ $46,2 \ km / h$}
\loigiai{
\begin{itemchoice}
\itemch 
\itemch 
\itemch 
\itemch 
\end{itemchoice}
}
\end{ex}
\begin{ex}
\immini[thm]{Lúc $7\ h$, một chiếc ca nô khởi hành từ $A$, di chuyển đến $B$ cách $A$ $12\ km$ dọc theo bờ một con sông. Coi chuyển động của ca nô là thẳng, tốc độ của dòng nước và tốc độ của ca nô đối với dòng nước không đổi. Chọn gốc thời gian lúc $7\ h$. Đồ thị độ dịch chuyển - thời gian của ca nô được cho như hình dưới.}{\begin{tikzpicture}[line cap=round, line join=round, scale=1, font=\footnotesize, >=stealth,draw=Mapcolor]
\path
(0,0)coordinate(O)
(4.5,0)coordinate(x)
(0,4)coordinate(y)

;
\draw  [dashed,help  lines, xstep=0.5cm,ystep=0.75cm]  (0,0)  grid  (4,4);
\draw [->](O)--(x);
\draw [->](O)--(y);

\node at (0,4.5) {$d\ (km)$};
\node at (5,0) {$t\ (h)$};

\foreach  \x  in  {0,1,2,...,4}
{\pgfmathtruncatemacro{\label}{\x};
\node[below]
(\label) at (\x,0) {\label};
\draw[fill=white](\x,0)circle(0.5pt);}

\foreach  \y  in  {1,2,...,4}
{\pgfmathtruncatemacro{\label}{3*\y};
\node[left]
(\label) at (0,.75*\y) {\label};
\draw[fill=white](0,.75*\y)circle(0.5pt);}


\draw[dashed] (0,3) -- (1,3) -- (1,0);
\draw[very thick,Mapcolor] (0,0) -- (1,3) -- (3,0);
\end{tikzpicture}}
\choiceTF[t]
{\True Ca nô trở lại bến $A$ lúc $10\ h$}
{Tốc độ trung bình của ca nô trong toàn bộ quá trình chuyển động là $0$}
{\True Chiều nước chảy là từ $A$ đến $B$}
{\True Nếu một chiếc bè bắt đầu trôi từ $A$ đến $B$ cùng lúc ca nô khởi hành thì nó đến $B$ lúc $11\ h$}
\loigiai{
\begin{itemchoice}
\itemch Dựa vào đồ thị độ dịch chuyển - thời gian: Ca nô trở lại bến A khi độ dịch chuyển $d=0$, tại thời điểm. Vậy ca nô trở lại bến A lúc 10 giờ.
$\rightarrow$ Đúng
\itemch Quãng đường chuyển động:
Tốc độ trung bình của ca nô trong cả chuyển động:
$\rightarrow$ Sai
\itemch Tốc độ của ca nô (so với bờ sông) khi đi theo chiều từ A đến B:
Tốc độ của ca nô (so với bờ sông) khi đi theo chiều từ B đến A:
Vì nên khi đi từ $A \rightarrow B$ thì ca nô đi xuôi dòng. Hay nước chảy từ $A \rightarrow B$.
$\rightarrow$ Đúng
\itemch Gọi vận tốc của ca nô đối với dòng nước là; gọi vận tốc của dòng nước đối với bờ sông là.
Ta có:
Một chiếc bè nếu trôi theo dòng nước thì vận tốc của nó chính là vận tốc của dòng nước.
Bè sẽ trôi từ bến A đến bến B trong: (kể từ thời điểm 7 giờ).
Nên bè sẽ đến bến B lúc 11 giờ.
\end{itemchoice}
}
\end{ex}
\begin{ex}
\immini[thm]{Ba chất điểm $(I)$, $(II)$, $(III)$ chuyển động giữa hai điểm $A$ và $B$ cách nhau $40\ m$. Đồ thị tọa độ - thời gian của các vật được biểu diễn như hình vẽ.}{
\begin{tikzpicture}[line cap=round, line join=round, scale=1, font=\footnotesize, >=stealth,draw=Mapcolor]
\path
(0,0)coordinate(O)
(4.5,0)coordinate(x)
(0,4)coordinate(y)

;
\draw  [dashed,help  lines, xstep=0.5cm,ystep=0.75cm]  (0,0)  grid  (4,4);
\draw [->](O)--(x);
\draw [->](O)--(y);

\node at (0,4.5) {$x\ (m)$};
\node at (5,0) {$t\ (s)$};

\foreach  \x  in  {0,1,2,...,4}
{\pgfmathtruncatemacro{\label}{2*\x};
\node[below]
(\label) at (\x,0) {\label};
\draw[fill=white](\x,0)circle(0.5pt);}

\foreach  \y  in  {1,2,...,4}
{\pgfmathtruncatemacro{\label}{10*\y};
\node[left]
(\label) at (0,.75*\y) {\label};
\draw[fill=white](0,.75*\y)circle(0.5pt);}

\draw[very thick,Mapcolor] (0,0) -- (2,3);
\draw[very thick,Mapcolor] (0,3) -- (4,0);
\draw[very thick,Mapcolor] (0,0) .. controls (1.5,0.3) and (2.55,0.5) .. (3.5,2.25);
\node at (2,3.2) {(I)};
\node at (3.8,2.4) {(III)};
\node at (4.2,0.4) {(II)};
\end{tikzpicture}
}
\choiceTF[t]
{Các chất điểm $(I)$, $(II)$, $(III)$ đều chuyển động thẳng đều}
{Vận tốc chất điểm $(I)$ là $v_1=10 \ m/s$ và vận tốc chất điểm $(II)$ là $v_2=5 \ m/s$}
{Phương trình tọa độ - thời gian của chất điểm (I) là: $x_1=40+10$ ($d:\ m$; $t: s$; $0\ s\leq t \leq 8\ s$}
{\True Chất điểm $(I)$ và $(II)$ gặp nhau tại thời điểm $t=\dfrac{8}{3}\ s$}
\loigiai{
\begin{itemchoice}
\itemch 
\itemch 
\itemch 
\itemch 
\end{itemchoice}
}
\end{ex}


\subsubsection*{PHẦN 3. Thí sinh trả lời từ câu 1 đến câu 6.}
\setcounter{ex}{0}
\begin{ex}%DoanhNamDinh
Một ô tô chạy từ địa điểm A đến địa điểm B với tốc độ $40 \ km / h$, sau đó ô tô quay trở về A với tốc độ $60 \ km / h$. Giả sử ô tô luôn chuyển động thẳng đều. Tốc độ trung bình của ô tô trên cả đoạn đường đi và về là bao nhiêu $m/s$ (làm tròn kết quả đến chữ số hàng phần mười)?
\shortans[oly]{13,3}

\loigiai{
$v_{tb}=\dfrac{s}{t}=\dfrac{2. S}{\dfrac{S}{v_1}+\dfrac{S}{v_2}}=\dfrac{2}{\dfrac{1}{40}+\dfrac{1}{60}}=48 \ km / h=13,3 \ m/s$
}
\end{ex}

\begin{ex}
Cho hai ô tô cùng lúc khởi hành ngược chiều nhau từ hai điểm $A$, $B$ cách nhau $120\ km$. Xe chạy từ $A$ với tốc độ $v_{A}=60\ km / h$, xe chạy từ $B$ với tốc độ $v_{B}=40\ km / h$. Hai xe gặp nhau sau mấy giờ chuyển động (làm tròn kết quả đến chữ số hàng phần mười)?
\shortans[oly]{1,2}
\loigiai{}
\end{ex}
\begin{ex}
Cùng một thời điểm, chất điểm thứ nhất chuyển động thẳng đều đi qua điểm $A$ với vận tốc $5 \ m/s$, chất điểm thứ hai chuyển động thẳng đều đi qua điểm $B$ với vận tốc $3 \ m/s$. Hai điểm $A$ và $B$ nằm trên cùng một đường thẳng, $B$ cách $A$ $50\ m$. Hai chất điểm chuyển động cùng chiều theo hướng từ $A$ đến $B$. Quãng đường chất điểm thứ nhất đi được cho đến khi gặp chất điểm thứ hai là bao nhiêu mét (làm tròn kết quả đến chữ số hàng đơn vị)?
\shortans[oly]{125}
\loigiai{}
\end{ex}

\begin{ex}
\immini[thm]{Đồ thị vận tốc - thời gian của một vật chuyển động được biểu diễn như hình vẽ. Quãng đường vật đi được từ thời điểm $t=0$ đến thời điểm $t=60\ s$ là bao nhiêu $km$?}{
\begin{tikzpicture}[line cap=round, line join=round, scale=1, font=\footnotesize, >=stealth,draw=Mapcolor]
\path
(0,0)coordinate(O)
(4.5,0)coordinate(x)
(0,3)coordinate(y)

;
\draw [->](O)--(x);
\draw [->](O)--(y);

\node at (0,3.5) {$v\ (m/s)$};
\node at (5,0) {$t\ (s)$};

\foreach  \x  in  {2,6,8}
{\pgfmathtruncatemacro{\label}{10*\x};
\node[below]
(\label) at (.5*\x,0) {\label};
\draw[fill=white](.5*\x,0)circle(0.5pt);}

\foreach  \y  in  {0,2,4}
{\pgfmathtruncatemacro{\label}{\y};
\node[left]
(\label) at (0,.5*\y) {\label};
\draw[fill=white](0,.5*\y)circle(0.5pt);}

\draw[dashed] (0,2) -- (1,2) -- (1,0);
\draw[dashed] (3,2) -- (3,0);
\draw[very thick,Mapcolor] (0,1) -- (1,2) -- (3,2) -- (4,0);
\end{tikzpicture}}
\shortans[oly]{2,2}
\loigiai{}
\end{ex}






\begin{ex}%Tailieuchuan
Một chiếc xe đạp đang chuyển động trên một đoạn đường thẳng, với độ dịch chuyển tại các thời điểm khác nhau được cho đồ thị như hình vẽ. Tỉ số giữa tốc độ và vận tốc trong suốt quá trình chuyển động là bao nhiêu (làm tròn kết quả đến chữ số hàng phần trăm)?
\begin{center} 
\begin{tikzpicture}[line cap=round, line join=round, scale=1, font=\footnotesize, >=stealth,draw=Mapcolor]
\draw  [help  lines, xstep=0.5cm,ystep=0.75cm]  (-0.5,-3)  grid  (10.5,3.5);
\foreach  \x  in  {1,2,3}
{\pgfmathtruncatemacro{\label}{2*\x};
\node[below]
(\label) at (\x,0) {\label};}

\foreach  \x  in  {4,5,6,7,8,9,10}
{\pgfmathtruncatemacro{\label}{3*\x};
\node[below]
(\label) at (\x,0) {\label};}

\foreach  \y  in  {-3,-2,...,4}
{\pgfmathtruncatemacro{\label}{5*\y};
\node[left=20pt]
(\label) at (0,.75*\y) {\label};}

\draw[very thick,Mapcolor] (0,0) -- (3,1.8) -- (4,1.8) -- (6,2.8) -- (7,2.8) -- (8,-1.5) -- (9,1.8) -- (10,1.8);
\node at (0,0.4) {$0$};
\node at (3,2.2) {$12$};
\node at (4,2.2) {$12$};
\node at (6,3.2) {$18$};
\node at (7,3.2) {$18$};
\node at (8,-1.8) {$-10$};
\node at (9,2.2) {$12$};
\node at (10,2.2) {$12$};
\node at (1,1) {$4$};
\node at (2,1.6) {$8$};
\node at (5,2.6) {$15$};
\node at (5.2,4.4) {Độ dịch chuyển (m) theo thời gian};
\node at (5.2,-3.6) {Thời gian (s)};
\node[rotate=90] at (-1.6,0.2) {Độ dịch chuyển (m)};
\end{tikzpicture}
\end{center} 
\shortans[oly]{5,67}
\loigiai{
}
\end{ex}
\begin{ex}
\impicinpar{
Một người đang chuyển động theo hướng từ $A$ đến $C$ với tốc độ không đổi $v_2=\dfrac{20}{\sqrt{3}}\ km / h$ tại thời điểm người đó đến $A$ cách đường quốc lộ $BC$ một đoạn $AH=200\ m$ thì nhìn thấy một xe ô tô vừa đến $B$ cách $A$ một đoạn $d=500\ m$ đang chạy trên đường với tốc độ không đổi như hình vẽ. Biết người đó và xe ô tô đến $C$ cùng lúc, cho góc $\alpha=120^\circ$. Vận tốc $v_1$ có độ lớn là bao nhiêu $m/s$ (làm tròn kết quả đến chữ số hàng đơn vị)?}{\begin{tikzpicture}[line cap=round, line join=round, scale=1, font=\footnotesize, >=stealth,draw=Mapcolor]
\path
(5.5,0)coordinate(C)
(3,2.5)coordinate(A)
(-1.5,0)coordinate(B)
(3,0)coordinate(H)
;

\draw[Mapcolor] (B)--(C) (A)--(C);
\draw[dashed,Mapcolor] (A)--(B) (A)--(H);
\draw pic[draw,angle radius=0.]{right angle=C--H--A};
\draw[very thick,->,red] (-1.5,0) -- (-0.5,0);

\draw[very thick,->,red]  (3,2.5) -- (4,1.5);
\node at (-1,-0.5) {$\vec v_1$};
\node at (4,2) {$\vec v_2$};
\node at (0.5,1.5) {$d$};
\node at (3.5,1) {$h$};

\draw pic[draw,angle radius=0.8cm]{angle=H--B--A};
\draw pic[draw,angle radius=0.4cm]{angle=B--A--C};
\draw pic["$\beta$",angle radius=1.85cm]{angle=H--B--A};
\draw pic["$\alpha$",angle radius=1.15cm]{angle=B--A--C};
\foreach \x/\g in {A/90,B/180,C/0}
\draw[fill=white](\x) circle(0.05)+(\g:0.3)node{$\x$};
\end{tikzpicture}}
\shortans[oly]{25}
\end{ex}








